\centerline{\bf \Huge Производная многочлена}

\begin{flushright}
        \scriptsize{Эпизод 6. Схватка в Токио-3}
\end{flushright}

\textbf{Определение.}
Производной многочлена $P(x)=a_n x^n+\ldots+a_1 x+a_0$ назовём многочлен $P^{\prime}(x)=n a_n x^{n-1}+(n-1) a_{n-1} x^{n-2}+\ldots+a_1$.

\textbf{Теорема 1.}
Если $P^{\prime}(x)>0$ на интервале $(a, b)$, то $P(x)$ возрастает на интервале $(a, b)$. Если $P^{\prime}(x)<0$ на интервале $(a, b)$, то $P(x)$ убывает на интервале $(a, b)$.

\textbf{Теорема 2.}
Если $x_0$ --- точка локального максимума или минимума многочлена $P(x)$, то $P^{\prime}\left(x_0\right)=0$.

\problem{1}
Верно ли утверждение, обратное утверждению теоремы 2 ?

\problem{2}
У многочлена степени $n$ имеется $n$ различных действительных корней. Докажите, что их среднее арифметическое равно среднему арифметическому действительных корней производной многочлена.


\problem{3}
Пусть $P(x)$ и $Q(x)$ - произвольные многочлены. Докажите, что

$$
(P(x) \cdot Q(x))^{\prime}=P^{\prime}(x) \cdot Q(x)+P(x) \cdot Q^{\prime}(x) .
$$


\textit{Напоминание.} Говорят, что $x_0$ является корнем многочлена $P(x)$ кратности $k$, если $P(x)=\left(x-x_0\right)^k Q(x)$, где $Q(x)$ - многочлен, причём $Q\left(x_0\right) \neq 0$.

\problem{4}
Докажите, что многочлен $P(x)$ степени $n \geqslant 1$ имеет корень кратности $k>1$ тогда и только тогда, когда $P(x)$ и $P^{\prime}(x)$ имеют общий корень.

\problem{5}
Дан многочлен $P(x)=1+x+\frac{x^2}{2!}+\ldots+\frac{x^n}{n!}$.

(a) Докажите, что $P(x)$ не имеет кратных корней.

(b) Докажите, что $P(x)$ имеет не более одного действительного корня.

\problem{6}
Даны многочлены $f$ и $g$ степени $n$. Докажите, что функция

$$
f g^{(n)}-f^{\prime} g^{(n-1)}+\ldots(-1)^n f^{(n)} g
$$


является константой. ( $g^{(k)}-k$-ая производная многочлена $g$.)

\newpage

\hspace{3mm}

\problem{7}
(a) Докажите, что у многочлена $P(x)=a_n x^{k_n}+a_{n-1} x^{k_{n-1}}+\ldots+a_1 x^{k_1}+a_0$ не более $n$ положительных корней.

\noindent
(b) Многочлен $P(x)$ степени $n$ имеет $n$ различных действительных корней. Какое наибольшее число его коэффициентов может равняться нулю?